\documentclass[tikz, border=6pt]{standalone}
% Shared preamble of the standalone TikZ figures (compiled by scripts/build_figures.py).
\usepackage{amsmath, amssymb}
\usepackage{xeCJK}
\setCJKmainfont[BoldFont={Hiragino Sans W6}]{Hiragino Sans W4}
\usetikzlibrary{arrows.meta, positioning, decorations.pathreplacing, calc, patterns}
\definecolor{rod}{RGB}{70,110,170}
\definecolor{rodb}{RGB}{200,120,50}
\definecolor{okc}{RGB}{60,140,80}
\definecolor{ngc}{RGB}{190,70,60}
\definecolor{boxc}{RGB}{235,240,248}
\tikzset{
  >={Stealth[length=2.4mm]},
  every node/.style={font=\small},
  box/.style={draw, rounded corners=2pt, fill=boxc, align=flush center, inner sep=4pt},
  rodA/.style={draw, fill=rod!35},
  rodB/.style={draw, fill=rodb!35},
  note/.style={font=\footnotesize, align=flush center},
  expl/.style={font=\small, align=flush left},
  brace/.style={decorate, decoration={brace, amplitude=4pt}},
}

\begin{document}
\begin{tikzpicture}
  % (a) comparison
  \node[anchor=west, font=\bfseries] at (-0.6, 1.4) {(a) 比較 $a \preceq b$};
  \draw[dashed, gray] (0, -1.0) -- (0, 1.0);
  \node[note, gray, anchor=north] at (0, -1.0) {一端をそろえる};
  \draw[rodA] (0, 0.35) rectangle (3.0, 0.75) node[midway] {$a$};
  \draw[rodB] (0,-0.55) rectangle (4.2,-0.15) node[midway] {$b$};
  \draw[dashed, gray] (3.0, 0.35) -- (3.0, -0.75);
  \node[note, anchor=west] at (3.0, 0.55) {\ $a$ の他端は $b$ の他端を越えない};

  % (b) concatenation
  \begin{scope}[xshift=8.2cm]
    \node[anchor=west, font=\bfseries] at (-0.6, 1.4) {(b) 連結 $a \circ b$};
    \draw[rodA] (0, 0) rectangle (3.0, 0.4) node[midway] {$a$};
    \draw[rodB] (3.0, 0) rectangle (7.2, 0.4) node[midway] {$b$};
    \draw[brace, decoration={mirror}] (0,-0.1) -- (7.2,-0.1) node[midway, below=5pt] {$a \circ b$：端と端でつなぐ};
  \end{scope}

  % (c) scale
  \begin{scope}[yshift=-3.4cm]
    \node[anchor=west, font=\bfseries] at (-0.6, 1.3) {(c) 目盛り：単位の棒 $u$ を並べて挟む};
    \foreach \i in {0,...,4} {
      \draw[draw, fill=gray!20] (\i*1.3, 0.3) rectangle ({(\i+1)*1.3}, 0.7) node[midway] {$u$};
    }
    \draw[rodA] (0,-0.5) rectangle (5.7,-0.1) node[midway] {$r$};
    \draw[dashed, gray] (5.2, 0.8) -- (5.2,-0.7) node[below, note, black] {$4u$};
    \draw[dashed, gray] (6.5, 0.8) -- (6.5,-0.7) node[below, note, black] {$5u$};
    \node[expl, anchor=west, text width=6.4cm] at (7.4, 0.1)
      {$4u \preceq r \preceq 5u$ なので $4 \le \ell(r) \le 5$。\\
       $N$ 本の $r$ と $M$ 本の $u$ を比べれば、$\ell(r)$ は有理数 $M/N$ でいくらでも細かく挟める。\\
       これが目盛りの一意性の証明である。};
  \end{scope}
\end{tikzpicture}
\end{document}
